\documentclass[12pt,a4paper]{ctexart}
\usepackage{geometry}
\usepackage{fontspec}
\geometry{left=2.5cm,right=2.5cm,top=2.5cm,bottom=2.5cm}
\linespread{1.3}
\setmainfont{Times New Roman}
\setCJKmainfont{SimSun}[AutoFakeBold=2.5]
\begin{document}
\pagestyle{empty}
\begin{center}
	{\heiti\zihao{3} AI 工具使用详情}
\end{center}
\vspace{6pt}
{\heiti\zihao{-4} 一、所用 AI 工具名称与版本}
\par
本参赛队在竞赛过程中使用 DeepSeek V4.1 Flash（对话式大语言模型，网页版/API 版），未使用其他 AI 工具。
\par
{\heiti\zihao{-4} 二、具体使用目的和环节}
\par
（1）代码调试：定位程序报错、提示边界条件与退化情形、建议自测方法；不代写核心算法。
\par
（2）语言润色：对论文语句的通顺性与表述一致性提出修改建议，不改变数据与结论。
\par
（3）经典算法核对：查阅经典算法的标准形式与符号约定（Welzl 最小外接圆、Andrew 单调链、旋转卡壳、Huber 加权最小二乘等），用于与自写实现相互核对。
\par
{\heiti\zihao{-4} 三、主要提示方式与使用过程}
\par
提问均以参赛队自己的代码片段或公式片段为上下文，并限定其只指出错误与检验方法。典型示例：把自写的旋转卡壳函数片段交给模型，询问“当凸包存在共线点、直径由多条平行支撑线给出时，是否会漏掉直径端点”；据此在实现中补充了共线退化处理，并增加暴力枚举对照断言。另有示例：提交虚拟时间分解代码与若干局日志，询问“累计值与日志字段不一致时可能的记账错误位置”，据此发现一处重试请求未按幂等处理的问题。
\par
{\heiti\zihao{-4} 四、对 AI 输出的采纳、人工修改和核验情况}
\par
所有题面理解、建模思路、模型建立、算法设计、实验方案与结果分析均由参赛队独立完成。AI 输出仅用于定位错误、提示边界与润色语句，采纳前逐项人工核验：
\par
（1）代码类修改一律以独立对照复核：旋转卡壳与暴力枚举逐位一致（偏差 $1.137\times10^{-13}$ m）；离线复现器与全部演练日志的计时口径相对误差 0.00\%；问题三、四各局结果均以 100\% 清除率为验收条件。
\par
（2）公式与符号均与原始文献核对（Huber 1964、Wilson 1927、Jung 1901、Church \& ReVelle 1974 等），确认无误后写入论文。
\par
（3）被 AI 建议但实测不利的方案一律不予保留：就近补测（11 源 7724 s、失败 29 次）、把补测触发条件扩展到交会角病态（离线 150 局三种阈值均略差于基准）等，均作为“检验驱动”的负结果记录在论文中。
\par
（4）语言润色仅调整语句表达；所有数字、图表与结论以参赛队自算结果为准，润色前后数据一致。
\end{document}
